% Progetto di esempio — LaTeX Editor Offline
% Contenuto originale, liberamente modificabile.
\documentclass[11pt,a4paper]{article}
\usepackage[T1]{fontenc}
\usepackage[utf8]{inputenc}
\usepackage[italian]{babel}

\title{Progetto dimostrativo}
\author{LaTeX Editor Offline}
\date{\today}

\begin{document}
\maketitle

\begin{abstract}
  Progetto dimostrativo per verificare l'editor: apertura, modifica,
  compilazione e anteprima PDF.
\end{abstract}

\section{Introduzione}
Questo progetto serve a provare il flusso completo dell'editor:
aprire la cartella, modificare questo file, salvare con \texttt{Ctrl+S}
e lanciare la build. La citazione d'esempio è \cite{rossi2024}.

\section{Strutture di esempio}

\subsection{Elenco}
\begin{itemize}
  \item primo elemento;
  \item secondo elemento.
\end{itemize}

\subsection{Tabella}
\begin{table}[htbp]
  \centering
  \begin{tabular}{lrr}
    \hline
    Strumento & Prima colonna & Seconda colonna \\
    \hline
    Editor    & 1             & 2                 \\
    Compilatore & 3           & 4                 \\
    \hline
  \end{tabular}
  \caption{Tabella di esempio.}
  \label{tab:esempio}
\end{table}
La tabella~\ref{tab:esempio} mostra dati fittizi.

\subsection{Figura senza file esterni}
\begin{figure}[htbp]
  \centering
  \fbox{\rule{0.7\linewidth}{0.35\linewidth}}
  \caption{Segnaposto grafico che non richiede file immagine.}
  \label{fig:segnaposto}
\end{figure}

\subsection{Matematica}
\begin{equation}
  \label{eq:demo}
  \int_{0}^{1} x^{2} \, dx = \frac{1}{3}
\end{equation}
L'equazione~\eqref{eq:demo} è un riferimento interno.

\section{Conclusioni}
Se il PDF viene generato e mostrato nel pannello destro, tutto funziona.

\bibliographystyle{plain}
\bibliography{bibliography}

\end{document}
